\section{Desarrollo del trabajo por per\'iodos}
\label{sec:cronologia}

Esta secci\'on documenta \emph{todo} lo realizado, agrupado por per\'iodos de trabajo, e
incluye los caminos explorados y abandonados, que son parte del trabajo de ingenier\'ia
aunque no aparezcan en el resultado final. La fuente es la bit\'acora del proyecto, que
registra \textbf{66 jornadas de trabajo activo} y aproximadamente \textbf{230
\emph{commits}} entre el 25 de febrero y el 1 de agosto de 2026.

\begin{figure}[ht]
\centering
\includegraphics[width=\textwidth]{timeline_commits.png}
\caption{Cronolog\'ia del desarrollo reconstruida desde la bit\'acora y los hitos del
repositorio. La densidad de actividad contextualiza los per\'iodos narrados; no se usa
como m\'etrica de calidad ni de avance cient\'ifico.}
\label{fig:timeline-desarrollo}
\end{figure}

\subsection{Del 25 de febrero al 23 de marzo --- Construcci\'on del marco te\'orico}

En este per\'iodo se cre\'o el repositorio del proyecto y se construy\'o la base te\'orica y
bibliogr\'afica del m\'etodo. Se mont\'o un \emph{vault} de conocimiento estructurado en
Obsidian y se resumieron \'integramente los ocho cap\'itulos del texto de referencia de
Foti \emph{et al.}, produciendo alrededor de 3000 l\'ineas de material propio: desde los
fundamentos de mec\'anica del continuo y propagaci\'on de ondas hasta la inversi\'on y los
casos de estudio. Se desarroll\'o en detalle el Problema de Lamb (1904) y el tratamiento
del n\'umero de onda complejo en medios viscoel\'asticos, y se estableci\'o el criterio
cuantitativo del efecto de campo cercano, que luego condicionar\'ia la geometr\'ia de los
arreglos de medici\'on.

Se escribieron cinco \emph{scripts} MATLAB de verificaci\'on te\'orica y se incorpor\'o como
subm\'odulo el c\'odigo MATLAB de MASW disponible p\'ublicamente, como banco de contraste.
Paralelamente se construy\'o una base de datos bibliogr\'afica que creci\'o de 22 a
\textbf{65 art\'iculos} curados, cada uno con an\'alisis estructurado individual, junto con
34 PDFs. De ese relevamiento surgieron los trabajos fundacionales que orientan el
proyecto (Park \emph{et al.}~1999, Xia \emph{et al.}~1999, las gu\'ias
InterPACIFIC~2018) y los antecedentes regionales de validaci\'on cruzada MASW--SPT.

\subsection{Del 10 al 15 de abril --- Consolidaci\'on del conocimiento}

Se ejecut\'o un proceso asistido y supervisado de consolidaci\'on del \emph{vault}: se
enriquecieron los conceptos existentes, se ampli\'o el conjunto de conceptos at\'omicos
interconectados hasta superar los \textbf{140}, se integraron los art\'iculos de la base
bibliogr\'afica dentro de los cap\'itulos correspondientes, se agregaron ejemplos trabajados
por cap\'itulo y se realiz\'o una pasada de enlazado cruzado. Se embebieron adem\'as las
figuras del texto de referencia y se reorganiz\'o el repositorio de recursos gr\'aficos.

\nota{Este bloque se apoy\'o fuertemente en asistencia de IA bajo supervisi\'on. Conviene
decidir con el tutor cu\'anto explicitar en el documento entregado y con qu\'e encuadre
metodol\'ogico.}

\subsection{Del 17 al 30 de abril --- Transductor, primeras se\~nales y dos pivotes de dise\~no}

Aqu\'i comienza el trabajo de hardware. Se model\'o el ge\'ofono SM-24 en Simulink a partir
de su funci\'on de transferencia de segundo orden y se cre\'o el primer proyecto en PSoC
Creator. El problema de dise\~no qued\'o planteado de inmediato: c\'omo extender la banda
\'util por debajo de la frecuencia natural del transductor, donde reside la informaci\'on de
las capas profundas.

\textbf{Se explor\'o y descart\'o un primer camino.} Se dise\~n\'o y optimiz\'o un filtro pasivo
RC de entrada; a las 24 horas se elimin\'o en favor de un esquema de \textbf{FIR digital
con decimaci\'on en cadena}, m\'as flexible y libre de tolerancias de componentes.

Se ejecut\'o una sesi\'on intensiva de desarrollo de firmware que produjo la arquitectura de
banco de pruebas del PSoC, un FIR en formato Q1.15 cargable desde el anfitri\'on, la
primera arquitectura de red ESP32 sobre WiFi, y una suite de \textbf{125 pruebas
automatizadas}. Esa suite detect\'o defectos reales que de otro modo habr\'ian llegado a
campo: un desbordamiento en la conversi\'on Q1.15, un error en el estado del bloque de
filtro por canal y una divisi\'on por cero que produc\'ia NaN.

Se obtuvo la \textbf{primera se\~nal real del SM-24}. Para poder verla se desarroll\'o
\texttt{geophone\_scope.m}, un osciloscopio en MATLAB sobre puerto serie de 657 l\'ineas,
y m\'as tarde una versi\'on simplificada. Se implementaron dos compensadores de extensi\'on
de banda por caminos distintos ---uno anal\'itico y otro por b\'usqueda exhaustiva--- y se
analiz\'o un filtro VCVS como compensador activo, con su funci\'on de transferencia completa
y el an\'alisis de compromisos frente a la alternativa digital.

Se construy\'o una prueba de concepto de la cadena diferencial completa
(\texttt{ADC $\to$ DFB $\to$ FIR notch $\to$ UART}), con un \textbf{filtro FIR
\emph{notch} de \SI{50}{\hertz} de 465 coeficientes} dise\~nado con \texttt{firls} para
suprimir la interferencia de red el\'ectrica, y se realizaron 18 sesiones de captura.

\textbf{El segundo pivote} cerr\'o el per\'iodo: se abandon\'o la topolog\'ia Sallen-Key/VCVS y
se adopt\'o \textbf{MFB} (\emph{Multiple Feedback}), decisi\'on respaldada por un an\'alisis
comparativo cuantitativo, junto con la suma en corriente mediante amplificador de
transimpedancia. Se realizaron las \textbf{primeras capturas en exteriores}.

\subsection{Del 5 al 20 de mayo --- Caracterizaci\'on espectral y cadena anal\'ogica completa}

Se desarroll\'o \texttt{analisis\_espectral.m}, una herramienta de estimaci\'on de densidad
espectral de potencia por el m\'etodo de Welch y c\'alculo de relaci\'on se\~nal--ruido por
entidad de medici\'on, y con ella se realiz\'o la comparaci\'on experimental sistem\'atica
entre las topolog\'ias MFB y Sallen-Key sobre datos reales. Se dise\~n\'o un filtro FIR
multibanda de orden 800 y se acumularon 41 sesiones de captura.

Se logr\'o la \textbf{calibraci\'on manual del \emph{offset} diferencial} mediante ajuste de
las referencias VDAC de cada etapa ---el antecedente directo de la auto-calibraci\'on
autom\'atica desarrollada un mes despu\'es--- junto con un protocolo de configuraci\'on
PSoC$\to$MATLAB. El protocolo de transmisi\'on evolucion\'o a 4 bytes con tipo y canal
multiplexados, y se agreg\'o un multiplexor anal\'ogico para medir el modo com\'un.

Un hallazgo relevante de este per\'iodo: el modelo de amplificador operacional ideal no
reproduc\'ia el comportamiento medido del compensador en el extremo superior de la banda.
Se incorpor\'o un primer \textbf{modelo no ideal del operacional} ($A_0=\SI{90}{\decibel}$
y un par\'ametro de \SI{8}{\mega\hertz}). Una auditor\'ia posterior detect\'o que el
segundo valor se trat\'o como polo dominante aunque dimensionalmente corresponde a una
frecuencia de ganancia unitaria. La relaci\'on $f_p=f_t/A_0$ ya se corrigi\'o en el c\'odigo,
pero el ajuste debe recalcularse y contrastarse. Se atac\'o el dise\~no del
compensador con cuatro estrategias de optimizaci\'on en paralelo.

El per\'iodo cierra con \texttt{DiferencialToSingleEnded}: la \textbf{cadena anal\'ogica
completa} implementada en el PSoC, con la secuencia correcta de arranque de todos los
bloques anal\'ogicos.

\subsection{Del 24 de mayo al 10 de junio --- El sistema inal\'ambrico multi-nodo}

Se dise\~n\'o e implement\'o la red: arquitectura maestro--esclavo sobre \textbf{ESP-NOW},
con enlace al PSoC y comunicaci\'on con el anfitri\'on. Sigui\'o una depuraci\'on intensiva
sobre hardware real, que incorpor\'o un \emph{beacon} de diagn\'ostico, una sonda de arranque
con confirmaci\'on y medici\'on de tiempo de ida y vuelta, almacenamiento y reenv\'io, y
pulsos de depuraci\'on por hardware para poder observar la temporizaci\'on con osciloscopio.

El 26 de mayo se alcanz\'o el \textbf{primer funcionamiento de extremo a extremo}: con la
incorporaci\'on de la fase de preparaci\'on previa al disparo, el sistema transmiti\'o datos
reales del PSoC hasta el anfitri\'on por primera vez.

En paralelo se cre\'o el proyecto PSoC dedicado al \textbf{martillo instrumentado} para el
martillo de impacto con sensor de fuerza PCB~Piezotronics~086D20. Poco despu\'es se tom\'o una decisi\'on de arquitectura
de software importante: unificar el firmware de ambos tipos de nodo en un \textbf{\'unico
archivo} que detecta autom\'aticamente en compilaci\'on qu\'e hardware tiene debajo, evitando
la bifurcaci\'on del c\'odigo embebido.

Se realiz\'o una refactorizaci\'on mayor que migr\'o el enlace PSoC$\leftrightarrow$ESP32 de
\textbf{SPI a UART}, ampli\'o el contador de muestras a 16 bits e introdujo una m\'aquina de
estados en el PSoC. Luego se corrigieron tres problemas de radiofrecuencia y
temporizaci\'on que estabilizaron el sistema (frecuencia del \emph{beacon}, almacenamiento
y reenv\'io, y velocidad de la UART).

El per\'iodo cierra con dos hitos: la \textbf{interfaz web embebida} ---el maestro sirve
una aplicaci\'on de p\'agina \'unica sobre su propio punto de acceso, eliminando la necesidad
de llevar una computadora al campo--- y un an\'alisis formal de robustez de la
sincronizaci\'on que identific\'o un defecto de confirmaci\'on-antes-de-GPIO y estableci\'o el
presupuesto de \emph{jitter} entre nodos.

\subsection{Del 14 al 30 de junio --- Auto-calibraci\'on del \emph{front-end}}

\begin{figure}[ht]
\centering
\includegraphics[width=\textwidth]{autocalibration_signal_waterfall_composite.png}
\caption{Evidencia visual integrada del tramo junio--julio: trayectoria de
auto-calibraci\'on, captura temporal y visualizaci\'on tipo \emph{waterfall}. El montaje
re\'une resultados de etapas distintas para mostrar la progresi\'on de banco a campo.}
\label{fig:cronologia-autocal-campo}
\end{figure}

Este per\'iodo produjo el aporte metodol\'ogico central del proyecto, y es un buen ejemplo
de iteraci\'on de dise\~no porque \textbf{el algoritmo se reemplaz\'o dos veces}.

La primera implementaci\'on us\'o \textbf{b\'usqueda binaria} (bisecci\'on) sobre los c\'odigos
VDAC, y se valid\'o en hardware real. Sobre esa base se realizaron dos avances
estructurales: la adquisici\'on migr\'o de una ISR cl\'asica a una \textbf{cadena gobernada
por DMA con filtro FIR embebido}, dejando al procesador fuera del camino de datos; y la
calibraci\'on gan\'o \textbf{persistencia en EEPROM con CRC-16}, sobreviviendo a un
reinicio. La cadena DMA se generaliz\'o luego a un multiplexor de cuatro canales
(crudo/filtrado $\times$ individual/lote) y se activ\'o el \textbf{promediado por
hardware}, de modo que la calibraci\'on dej\'o de promediar muestra a muestra en la
interrupci\'on y pas\'o a consumir ventanas ya sumadas, con verificaci\'on de coherencia en
tiempo de compilaci\'on.

El 23 de junio se produjo el \textbf{reemplazo total del algoritmo}. Tras convivir con
bisecci\'on binaria, un servo lento y variantes de control PID, el sistema converg\'io en
un \'unico
\textbf{controlador PI} con prealimentaci\'on f\'isica y detector de enganche, unificado
para ambos tipos de nodo, y se document\'o la primera calibraci\'on exitosa sobre hardware
real.

El refinamiento posterior aport\'o: escala del error en c\'odigos de DAC, banda muerta
proporcional a la ganancia de cada etapa, \textbf{anti-windup direccional} ---el
integrador se congela contra el l\'imite del DAC \'unicamente en el sentido que empeorar\'ia
la saturaci\'on--- y manejo correcto de ganancias con signo. El controlador se generaliz\'o
a las cuatro etapas de la placa de ge\'ofono y se elimin\'o definitivamente el c\'odigo de
bisecci\'on heredado.

Dos hallazgos de este per\'iodo merecen menci\'on por su valor pr\'actico. Primero, se
descubri\'o que el filtro de adquisici\'on llevaba \emph{meses} cargando coeficientes de
pasa-banda en lugar del pasa-bajos que correspond\'ia. Segundo, un problema que aparentaba
ser de firmware result\'o ser \textbf{un pin sin soldar}. El per\'iodo cierra corrigiendo
cuatro defectos de campo documentados y agregando detecci\'on autom\'atica del tipo de nodo,
propagada desde el PSoC hasta la interfaz web.

\subsection{Del 1 al 12 de julio --- Determinismo, almacenamiento y procesamiento MASW}

El modelo de captura pas\'o de estar gobernado por temporizador a basarse en
\textbf{confirmaci\'on real de fin de captura por nodo}, y la EEPROM de calibraci\'on se
redise\~n\'o para mantener \textbf{un registro independiente por c\'odigo de ganancia}.

El cambio de arquitectura m\'as importante fue el nacimiento de \textbf{\emph{superMaquina}}:
la l\'ogica de captura migr\'o de C a una m\'aquina de estados descrita en \textbf{Verilog} e
instanciada en los bloques UDB del PSoC, que decide en hardware qu\'e solicitud de DMA se
deja pasar y cu\'ando arranca y termina una captura. La motivaci\'on es el determinismo: en
un sistema multicanal el \emph{jitter} temporal se traduce directamente en error de
velocidad de fase. Se agreg\'o adem\'as un conjunto completo de comandos por USB para operar
el \emph{pipeline} paso a paso desde terminal, sin depender de la web.

Se consolid\'o la pol\'itica de eventos ``determinismo primero'': los eventos de fin de
cuenta de los temporizadores dejaron de generar interrupci\'on y pasaron a consumirse por
sondeo, se dedic\'o un temporizador al \emph{watchdog} de captura, y se agreg\'o un
\textbf{cortacircuitos anti-tormenta de interrupciones} ---una condici\'on efectivamente
observada en placa mediante depuraci\'on por pyOCD--- con auto-recuperaci\'on. El ADC pas\'o
de dos a \textbf{cuatro configuraciones de rango de entrada}, todas validadas en
hardware, y la interfaz web sum\'o indicador de intensidad de se\~nal del enlace.

En paralelo arranc\'o el \emph{pipeline} de procesamiento en Python, que sum\'o cerca de
\textbf{4600 l\'ineas} de una vez: m\'odulos de an\'alisis de dispersi\'on e inversi\'on MASW y
una aplicaci\'on de revisi\'on de datos de campo. Esta \'ultima recibi\'o doce fases de pulido
(ordenamiento por amplitud pico a pico, banco de filtros, alineaci\'on entre tandas,
\emph{waterfall} y $f$--$k$, selecci\'on de curvas) y luego la capacidad de definir grupos
de dispersi\'on independientes con combinaci\'on ponderada de im\'agenes, base de la
inversi\'on multimodal. Se integr\'o \texttt{maswavespy} como subm\'odulo y se corri\'o la
primera deconvoluci\'on sobre una captura real de campo.

El per\'iodo cierra con la \textbf{auditor\'ia pre-campo} descrita en detalle en la
Secci\'on~\ref{sec:auditoria}, que encontr\'o y corrigi\'o varios defectos cr\'iticos y emiti\'o
el veredicto de sistema listo para campo.

\subsection{Del 15 al 29 de julio --- Comunicaci\'on cient\'ifica y reorganizaci\'on}

Se redact\'o el art\'iculo sobre la auto-calibraci\'on para el congreso URUCON~2026,
incorporando datos experimentales reales de dos placas con barrido completo de ganancia.
El art\'iculo se someti\'o a un panel de revisi\'on editorial simulada y se respondi\'o a las
observaciones: siete referencias nuevas, tabla de resultados ampliada a las cuatro etapas
en lugar de solo la \'ultima, ecuaciones de banda muerta con justificaci\'on expl\'icita, y
una correcci\'on de honestidad cient\'ifica sobre qu\'e representa realmente la condici\'on de
referencia. Tras una reuni\'on con el tutor se comprimi\'o de seis a cinco p\'aginas.

En el firmware se salda una deuda de nomenclatura renombrando por completo la etapa
sumadora, y se hace configurable la banda muerta por etapa. En MATLAB se establece un
protocolo formal de barridos de la cadena anal\'ogica con \emph{script} de an\'alisis y
datos reales, que permiti\'o contrastar la respuesta medida contra el modelo del circuito
con valores de componentes reales.

Se produjo tambi\'en una \textbf{ruptura arquitect\'onica} del repositorio: el historial
monol\'itico se preserv\'o y el proyecto se reinici\'o como superproyecto con siete
subm\'odulos independientes, reorganizados poco despu\'es \textbf{por prop\'osito y no por
lenguaje} (firmware, interfaces, c\'alculos y modelados).

Finalmente se port\'o el servidor Python de una implementaci\'on artesanal a
\textbf{FastAPI}, con un \emph{gate} automatizado como condici\'on de avance, y se
migr\'o la aplicaci\'on de escritorio PyQt de revisi\'on de campo a la web hasta alcanzar
paridad funcional, incorporando cuarentena reversible en lugar de borrado f\'isico, cola
de trabajos de an\'alisis persistente, exportaciones reproducibles y un \emph{gate} de
\textbf{58 verificaciones}.

\subsection{1 de agosto --- Hacia el hardware definitivo y la fuente s\'ismica}

Se migr\'o el conexionado PSoC$\leftrightarrow$ESP32 a una \textbf{placa portadora
intermedia} con asignaci\'on de pines definitiva verificada por \emph{script}, primer paso
hacia una PCB propia. Este trabajo requiri\'o reconstruir por ingenier\'ia inversa los
valores de componentes del esquem\'atico anal\'ogico, dado que el archivo del dise\~no es un
binario propietario sin herramientas de lectura p\'ublicas; la reconstrucci\'on se realiz\'o
contrastando el binario con la asignaci\'on generada por el ajustador y las cabeceras de
pines, declarando expl\'icitamente el nivel de confianza de cada valor recuperado. El
resultado es que el dise\~no anal\'ogico dej\'o de estar atrapado en un formato cerrado.

En paralelo arranc\'o el dise\~no del \textbf{martillo de leva}, la fuente s\'ismica
repetible, atacado simult\'aneamente con un modelo anal\'itico 2D en Python y un modelo 3D
en Simscape Multibody derivado del CAD real. Ambos modelos detectaron errores propios ---
una interferencia axial entre la leva y el brazo, y una violaci\'on de conservaci\'on de
energ\'ia--- que est\'an pendientes de resoluci\'on.

\subsection{S\'intesis: caminos explorados y descartados}

Por su valor metodol\'ogico se listan expl\'icitamente las alternativas evaluadas y
abandonadas, con el motivo del descarte.

\begin{table}[h]
\centering\footnotesize
\caption{Alternativas de dise\~no evaluadas y descartadas durante el proyecto.}
\label{tab:descartes}
\begin{tabularx}{\textwidth}{@{}l X l@{}}
\toprule
\textbf{Alternativa} & \textbf{Motivo del descarte} & \textbf{Per\'iodo} \\
\midrule
Filtro pasivo RC de entrada & Reemplazado por FIR digital: tolerancias, deriva t\'ermica, fase no lineal y disponibilidad de DSP & abr. \\
Topolog\'ia Sallen-Key / VCVS & An\'alisis comparativo favorece MFB (comportamiento con ganancia y sensibilidad a componentes) & abr. \\
Topolog\'ia VCVS de la compensaci\'on & Sensibilidad a tolerancias frente a la MFB; la compensaci\'on anal\'ogica activa se mantuvo & abr.--may. \\
Filtro adaptativo LMS para \emph{notch} & Transitorio de convergencia y dependencia de la se\~nal; reemplazado por regresi\'on arm\'onica por lote & jun. \\
\emph{Notch} arm\'onico en el cliente web & Retardo de grupo y eliminaci\'on indiscriminada dentro de la banda MASW; el crudo se trata en post-proceso & jun. \\
Enlace SPI PSoC$\leftrightarrow$ESP32 & Migrado primero a UART; la versi\'on actual usa UART en bajada e I\textsuperscript{2}C en subida & may.--ago. \\
Calibraci\'on por bisecci\'on binaria & Reemplazada por PI unificado; no manejaba bien ruido ni signo de ganancia & jun. \\
Servo lento y variantes PID & Se descarta el t\'ermino derivativo: amplifica el ruido de medici\'on del ADC sin aportar en una planta cuasi-est\'atica; el controlador final es PI & jun. \\
Captura gobernada por temporizador & Reemplazada por confirmaci\'on real de fin por nodo & jul. \\
L\'ogica de captura en C sobre registro & Reemplazada por m\'aquina de estados en Verilog (determinismo) & jul. \\
Servidor HTTP artesanal & Migrado a FastAPI & jul. \\
Aplicaci\'on de escritorio PyQt & Portada a la web para unificar la interfaz de operador & jul. \\
Repositorio monol\'itico & Superproyecto con subm\'odulos organizados por prop\'osito & jul. \\
\bottomrule
\end{tabularx}
\end{table}

\nota{Esta secci\'on completa no entra en las 15 p\'aginas. Propuesta: para la entrega
formal reducirla a \emph{una p\'agina}, con un p\'arrafo por per\'iodo, m\'as la
Tabla~\ref{tab:descartes} (que es la que mejor muestra criterio de ingenier\'ia). El
texto \'integro se conserva para el libro de tesis de la Defensa Final.}
